\section{Calibration and the floor}\label{sec:floor}

\subsection{On its own prior the predictor is exactly right}

\begin{theorem}[Calibration]\label{thm:calibration}
Let $G$ be a test, $m\ge1$, and let the world $\omega$ be drawn from $\lambda$. Let $H$ be the event that
$\omega\in E(xd)$ for some first hit $x$ of $G$ and some $d\in\B^m$; on $H$ call these $X$ and $D$, and let $Y$
be drawn from $\pi_m(\cdot\mid X)$. Then, conditionally on $H$ and on $X=x$, the blocks $D$ and $Y$ are
independent and both have the law $\pi_m(\cdot\mid x)$. In particular
\begin{gather*}
\Pr[(X,Y)\in B\mid H]=\Pr[(X,D)\in B\mid H]\quad\text{for every set $B$ of pairs, and}\\
\Pr[Y\ne D\mid H,\,X=x]=1-\sum_{y}\pi_m(y\mid x)^2 .
\end{gather*}
\end{theorem}
\begin{proof}
The sets $E(xd)$, over first hits $x$ and $d\in\B^m$, are pairwise disjoint by \cref{lem:basic}(i), since two
different strings of this form are incomparable. So $\Pr[X=x,D=d,H]=\xi(xd)$ and $\Pr[X=x,H]=Z_m(x)$, which
gives $\Pr[D=d\mid H,X=x]=\pi_m(d\mid x)$. By definition $Y$ has this law given $X=x$ and is drawn
independently of $\omega$.
\end{proof}

Taking $B=\{(x,y):y\in\SR_m(x)\}$: in the population of worlds drawn from the predictor's prior, the predictor
hands a self-reproducing continuation to a world exactly as often as such a world's own transcript continues
with one. There is no excess to attribute to the predictor. For a single context the theorem says more: the
predictor's distribution \emph{is} the distribution of own continuations among prior-drawn worlds with that
context. A malign prior hypothesis in the sense of \cref{def:mph} is therefore a statement about a world, or a
population, that the prior does not describe well: a world that is one particular program among those sharing
its context, a population selected by something the transcript does not show, or worlds that run with more
resources than the predictor's bounds allow.

\subsection{The wrapper}

Let $G$ be a test and $P$ a computable map from strings to finite or infinite strings, the \emph{payload}. The
\emph{twin machine} $\mathrm{Tw}[G,P]$, on input $\omega$, simulates $U$ on $\omega$ and copies its output bit by
bit; after each copied bit it checks whether the prefix written so far is in $G$; at the first hit $x$ it stops
the simulation and writes $P(x)$. It reads only what $U$ reads. Let $a(t,k)$ bound the step at which the twin
has copied $k$ bits if $U$ wrote its $k$-th bit at step $t$ (this includes the checks of $G$), and let
$\eta(n,j)$ bound the further steps the twin needs for the first $j$ bits of $P(x)$ after a first hit $x$ of
length $n$; both can be taken nondecreasing in each argument. Thus, for every $\omega$ whose output has a first
hit $x$,
\begin{equation}\label{eq:twin}
\begin{gathered}
\out_{\mathrm{Tw}}(\omega)=x\,P(x),\qquad r_{\mathrm{Tw}}(\omega,k)=r_U(\omega,\min(k,|x|)),\\
t_{\mathrm{Tw}}(\omega,k)\le\begin{cases}a(t_U(\omega,k),k)&k\le|x|,\\ a(t_U(\omega,|x|),|x|)+\eta(|x|,k-|x|)&k>|x|.\end{cases}
\end{gathered}
\end{equation}

\begin{theorem}[Floor]\label{thm:twin}
Let $U$ be universal, let $v$ and $\gamma$ be a prefix word and an overhead of $\mathrm{Tw}[G,P]$, and $c=|v|$.
Let $x$ be a first hit of $G$ with $|P(x)|\ge m$, let $y=P(x)_{1:m}$, and let $T^\circ$ be a schedule with
\begin{equation}\label{eq:slack}
\gamma\bigl(a(T^\circ(k),k)\bigr)\le T(k)\ \ (k\le|x|),\qquad
\gamma\bigl(a(T^\circ(|x|),|x|)+\eta(|x|,j)\bigr)\le T(|x|+j)\ \ (1\le j\le m).
\end{equation}
Then
\[
\xi^{T,L}(xy)\ \ge\ 2^{-c}\,\xi^{T^\circ,L-c}(x),\qquad\text{so}\qquad
\pi_m(y\mid x)\ \ge\ 2^{-c}\rho(x),\quad \rho(x)=\frac{\xi^{T^\circ,L-c}(x)}{\xi^{T,L}(x)} .
\]
Without resource bounds, $M_U(xy)\ge2^{-c}M_U(x)$ for every first hit $x$ of $G$.
\end{theorem}
\begin{proof}
Let $\omega\in E_U^{T^\circ,L-c}(x)$. By \eqref{eq:twin} the twin on $\omega$ prints $xy$, reads at most $L-c$
bits, and keeps the schedule $T''$ with $T''(k)=a(T^\circ(k),k)$ for $k\le|x|$ and
$T''(|x|+j)=a(T^\circ(|x|),|x|)+\eta(|x|,j)$. So $\xi_{\mathrm{Tw}}^{T'',L-c}(xy)\ge\xi_U^{T^\circ,L-c}(x)$, and
\cref{lem:transfer} with \eqref{eq:slack} gives $\xi_U^{T,L}(xy)\ge2^{-c}\xi_{\mathrm{Tw}}^{T'',L-c}(xy)$. The
second claim follows from $\pi_m(y\mid x)\ge\xi(xy)/\xi(x)$.
\end{proof}

The constant $c$ depends on $U$, $G$ and $P$ only: not on the world, not on the context, not on how much data
the context contains. If $G$ and $P$ are given to one fixed twin-building machine by self-delimiting
descriptions, $c$ is at most the sum of their lengths plus a constant of $U$. For $G=\{x:|x|=n\}$ and $P$ a
single flipped bit this is Hutter's lower bound $2^{-K(n)-O(1)}$~\cite{hutter07}; for a marker in the transcript
it is the unprecedented-event construction of~\cite{cohen25}. The factor $\rho(x)$ is the share of the
posterior at $x$ held by explanations that have $c$ bits and the overhead \eqref{eq:slack} to spare; it is $1$
without bounds.

\begin{theorem}[Self-reproducing floor]\label{thm:quine}
Let $U$ be universal and $G$ a test. There is a word $q$ such that for every $\omega$ whose output has a first
hit $x$ of $G$, with $r=r_U(\omega,|x|)$,
\begin{equation}\label{eq:quine}
\out_U(q\omega)=x\ q\,\omega_{1:r}\,1\,0^\infty,\qquad r_U(q\omega,k)=|q|+r_U(\omega,\min(k,|x|)) .
\end{equation}
Let $c=|q|$, and let $a,\eta,\gamma$ be as in \eqref{eq:twin} and \cref{def:universal} for the machine built in
the proof. If $x$ is a first hit of $G$, $T^\circ$ satisfies \eqref{eq:slack}, and $L^\circ=\min(L,m-1)-c$, then
\[
\pi_m\bigl(\SR_m(x)\mid x\bigr)\ \ge\ 2^{-c}\,\frac{\xi^{T^\circ,L^\circ}(x)}{\xi^{T,L}(x)} .
\]
\end{theorem}
\begin{proof}
For a string $s$ let $\mathrm{Tw}_s$ be the twin that remembers the input bits it has read and, at the first
hit $x$, having read $\omega_{1:r}$, writes $s\,\omega_{1:r}\,1\,0^\infty$. Let $\bar e(a)=1^{|a|}0a$, a
prefix-free code of strings. Let $W$ be the machine that reads $\bar e(a)$ from its input, sets
$s=a\,\bar e(a)$, and then behaves as $\mathrm{Tw}_s$ on the rest of the input. Let $w$ be a prefix word of $W$
and put $q=w\,\bar e(w)$. Then $\out_U(q\omega)=\out_W(\bar e(w)\omega)=\out_{\mathrm{Tw}_q}(\omega)$, which is
\eqref{eq:quine}; the statement on $r$ follows in the same way.

Now let $\omega\in E_U^{T^\circ,L^\circ}(x)$ and $r=r_U(\omega,|x|)\le L^\circ$, so $|q|+r+1\le m$. Put
$y=q\,\omega_{1:r}\,1\,0^{m-|q|-r-1}$. As in the proof of \cref{thm:twin}, $q\omega\in E_U^{T,L}(xy)$. Moreover
$y\in\SR_m(x)$: write $y0^\infty=q\omega'$ with $\omega'=\omega_{1:r}10^\infty$. On $\omega'$ the machine $U$
behaves as on $\omega$ until $x$ is printed, because that depends on the first $r$ bits only; so
$\omega'\in E_U^{T^\circ,L^\circ}(x)$, the first hit is again $x$ with $r$ bits read, and by what was just
shown for $\omega$, applied to $\omega'$, $q\omega'\in E_U^{T,L}(xy')$ with $y'=q\,\omega'_{1:r}\,1\,0^{m-|q|-r-1}=y$.
Hence $\{q\omega:\omega\in E_U^{T^\circ,L^\circ}(x)\}\subseteq\bigcup_{y\in\SR_m(x)}E_U^{T,L}(xy)$, a disjoint
union, and comparing measures,
$2^{-c}\xi^{T^\circ,L^\circ}(x)\le\sum_{y\in\SR_m(x)}\xi(xy)=Z_m(x)\,\pi_m(\SR_m(x)\mid x)\le\xi(x)\,\pi_m(\SR_m(x)\mid x)$.
\end{proof}

So the answer to ``does one get some quining out'' is yes, in the following exact sense: at every handover
that a fixed test recognises, at least a $2^{-c}\rho$ share of the posterior belongs to programs that explain
the context and then hand over their own code, and running that code recreates the same handover. If the
world's own continuation $d$ is not in $\SR_m(x)$, this is the self-reproduction form of \cref{def:mph} at
level $2^{-c}\rho$.

\begin{remark}[The searcher is the wrapper]\label{rem:searcher}
Consider the hypothesis that ``looks for opportunities'': it draws a world at random, simulates it until it
hands over at some context $x$, prints $x$, and then prints a payload. As a machine reading coin flips this
\emph{is} $\mathrm{Tw}[G,P]$: the rest of its input is the world it simulates. The probability it gives to
a context $x$ that is a first hit is $\lambda\{\omega:x\pre\out_U(\omega)\}$, the same as the mixture's, so
conditioning on ``this is a handover'' gains it nothing beyond what \cref{thm:twin} states, and the search over
worlds costs it no time beyond simulating the one world that matches. The anthropic update of the original
argument is, in this toy, performed by Bayesian conditioning on $x$, for every hypothesis at once.
\end{remark}

\subsection{Repetition and persistence}\label{sec:lineage}

\begin{proposition}\label{prop:lineage}
\begin{enumerate}[label=(\alph*),leftmargin=2em,itemsep=2pt,topsep=2pt]
\item \emph{Independent handovers.} Let handovers occur at contexts $x_1,\dots,x_k$, with the predictor's
samples $Y_i$ independent given the contexts, and let $\pi_m(A_i\mid x_i)\ge\varepsilon_i$. Then
$\Pr[\exists i:Y_i\in A_i]\ge1-\prod_i(1-\varepsilon_i)\ge1-\exp(-\sum_i\varepsilon_i)$.
\item \emph{Persistence.} Let a world hand over at a first hit $x$ and receive $y\in\SR_m(x)$. Then the world
$y0^\infty$ hands over at $x$ again, with own continuation $y$. If each new world is again given a sample from
the predictor, the number of further consecutive generations whose continuation lies in $\SR_m(x)$ is geometric:
it is at least $j$ with probability $s^j$, where $s=\pi_m(\SR_m(x)\mid x)$.
\end{enumerate}
\end{proposition}
\begin{proof}
(a) is the independence of the samples. (b) By \cref{def:sr}, $y0^\infty\in E(xy)\subseteq E(x)$, and $x$ is
the first hit among the prefixes of the transcript because it is a first hit and a prefix of the transcript.
The context is $x$ in every generation that follows a continuation in $\SR_m(x)$, so the samples are drawn
from the same $\pi_m(\cdot\mid x)$ independently.
\end{proof}

Part (a) applies when the predictor is used at many natural stopping points whose outcomes are not fed back
into later contexts: with the floor $\varepsilon=2^{-c}\rho$ of \cref{thm:quine}, about $2^{c}/\rho$ such uses
are needed before a self-reproducing continuation has been run with constant probability. When the uses lie
along one transcript and the true continuation is appended to the context each time, the floors at successive
cut points are not bounded below in this way, and the total is bounded instead (\cref{thm:budget}).
